{{دیوان
| عنوان = سر اکبر حیدری، صدرِ اعظم حیدر آباد دکن کے نام
| ماخذ = ویکی ماخذ
| ماخذ_لنک = https://ur.wikisource.org/wiki/%D8%B3%D8%B1_%D8%A7%DA%A9%D8%A8%D8%B1_%D8%AD%DB%8C%D8%AF%D8%B1%DB%8C%D8%8C_%D8%B5%D8%AF%D8%B1%D9%90_%D8%A7%D8%B9%D8%B8%D9%85_%D8%AD%DB%8C%D8%AF%D8%B1_%D8%A2%D8%A8%D8%A7%D8%AF_%D8%AF%DA%A9%D9%86_%DA%A9%DB%92_%D9%86%D8%A7%D9%85
| مقابلہ = اقبال سائبر لائبریری (اقبال اکادمی پاکستان)
| مقابلہ_لنک = https://www.iqbalcyberlibrary.net
}}
’یوم اقبالؔ‘ کے موقع پر توشہ خانۂ حضور نظام کی طرف سے، جو صاحب صدر اعظم کے ماتحت ہے ایک ہزار روپے کا چیک بطور تواضع وصول ہونے پر

<poem>
تھا یہ اللہ کا فرماں کہ شکوہِ پرویز
دو قلندر کو کہ ہیں اس میں ملوکانہ صفات

مجھ سے فرمایا کہ لے، اور شہنشاہی کر
حُسنِ تدبیر سے دے آنی و فانی کو ثبات

مَیں تو اس بارِ امانت کو اُٹھاتا سرِ دوش
کامِ درویش میں ہر تلخ ہے مانندِ نبات

غیرتِ فقر مگر کر نہ سکی اس کو قبول
جب کہا اُس نے یہ ہے میری خدائی کی زکات!
</poem>
